% Ad Quintum Fratrem 2.9 --- headnote
% ad-quintum-fratrem-02-09, early February 54 BC, Rome

\textit{Cicero to his brother Quintus, written at Rome at the
beginning of February 54 BC --- Quintus by then had set out for
Caesar's camp in Gaul. The letter is brief, conversational, and
proverbial. Quintus's writing-tablets have demanded a letter; the
matter itself, Cicero says, gives him no subject, but conversation
between brothers is allowed to ramble. The verb \textit{alucinari}
(``to ramble, to babble at random'') gives the licence.}

\textit{Section 2 carries the famous pun. The Tenedians of the
island near Troy had asked the Senate for some restoration of
liberty; the Senate refused. ``Cut off by a Tenedian axe'' ---
\textit{securi Tenedia} --- is a proverb (the Tenedian axe was
notoriously sharp), turned to the present case where the
applicants were themselves Tenedian. Cicero notes who voted with
him: M. Bibulus, M. Calidius, M. Favonius --- all four men of the
optimate front. The Magnesians from Sipylum, by contrast, sang
Quintus's praises for blocking L. Sestius Pansa.}

\textit{Section 3 is the famous Cicero on Lucretius. ``Many flashes
of genius, yet much craft as well'' --- the Latin is
\textit{multis luminibus ingeni, multae tamen artis} (text
disputed; some editors read \textit{non multae tamen artis},
``yet of no great craft''). The reading on either side is brief
and Quintus had already sent his own appraisal. ``But ---- when
you come.'' That, like much of the brothers' literary criticism,
will be in person. The closing joke on Sallustius's
\textit{Empedoclea} (a tedious cosmological poem in Empedoclean
hexameters by Cn. Sallustius, friend of Cicero, not the historian)
is the brothers' shared playfulness on each other's reading.}
